\documentclass[11pt]{article}

\usepackage[margin=1in]{geometry}
\usepackage[T1]{fontenc}
\usepackage[utf8]{inputenc}
\usepackage{microtype}
\usepackage{booktabs}
\usepackage{amsmath}
\usepackage{amssymb}
\usepackage{graphicx}
\usepackage{xcolor}
\usepackage{hyperref}
\usepackage[numbers,sort&compress]{natbib}

\hypersetup{
  colorlinks=true,
  linkcolor=black,
  citecolor=blue!50!black,
  urlcolor=blue!50!black
}

\title{Brain-Aligned Memorability Signals Improve Video Generation Selection}
\author{
  Jawaun Brown\\
  Independent Research Project\\
  \texttt{jawaun@example.com}
}
\date{Draft regenerated June 1, 2026}

\begin{document}
\maketitle

\begin{abstract}
Generated videos can now be sampled cheaply, but creators still need to decide
which candidate is worth showing to humans. Standard evaluation metrics target
visual quality, prompt alignment, or generic preference; they do not directly ask
what a viewer will remember. We study whether a brain-aligned video
representation can supply a practical memorability selector for generated
videos. Using 1,022 analyzable clips from the 1,102-clip BOLD Moments dataset,
we learn a supervised memorability readout in TRIBE-predicted cortical-response
features and compare it with self-supervised video features, prompt-preservation
signals, and generation-time selector policies. TRIBE predicts held-out BOLD
Moments memorability at Spearman $\rho=0.403\pm0.061$, comparable to the V-JEPA
baseline at $\rho=0.395\pm0.037$. On a Wan2.2 candidate pool, TRIBE/BMD and
V-JEPA selectors overlap only partially, creating a useful adjudication setting
for blinded human evaluation. The current paper therefore makes a narrow claim:
brain-aligned features expose a compact, auditable memorability signal and define
a testable generated-video selection workflow. It does not yet claim that
TRIBE-selected generated videos are more memorable to humans, nor that direct
generator steering is solved.
\end{abstract}

\section{Introduction}

Video generation systems increasingly make it cheap to produce many plausible
variants for the same prompt. This shifts a practical bottleneck from generation
to selection: which candidate should a creator, researcher, or product system
show to people? Existing automatic evaluations emphasize text-video alignment,
image fidelity, temporal consistency, aesthetic quality, or broad preference.
These are necessary but incomplete when the target property is cognitive:
whether a viewer will remember the clip.

Memorability is a useful first target because it is empirically measurable and
partly stable across observers \citep{isola2011memorability,newman2020memento}.
BOLD Moments further links short naturalistic videos, memorability metadata, and
fMRI responses \citep{lahner2024boldmoments}. We use this setting to ask whether
a brain-aligned representation can provide a practical selection signal for
generated videos.

The intended contribution is deliberately narrow. We do not claim that TRIBE is
a human preference oracle, that memorability equals engagement, or that a learned
linear vector is the essence of memory. We treat TRIBE, V-JEPA, CLIP, and human
behavior as different representation frames over the same stimulus space. The
scientific question is whether an ordering learned in one frame transfers to
independent human judgment.

\section{Related Work}

\paragraph{Memorability.}
Image memorability work established that memorability is not merely private
taste; some images are consistently remembered better than others
\citep{isola2011memorability}. Video memorability extends the problem to dynamic
events and temporal decay, requiring both visual and semantic representations
\citep{newman2020memento,cohendet2019videomem}. MediaEval's video
memorability tasks further show that the benchmark ecosystem distinguishes
short-term memorability, dataset generalization, and EEG-adjacent prediction
subtasks \citep{sweeney2022mediaeval}. Applied systems such as MindMem are
useful commercial-adjacent memorability predictors, but they should not be read
as neuroscience validation \citep{asgarian2025mindmem}.

\paragraph{Brain-aligned video representations.}
BOLD Moments provides short naturalistic videos, fMRI responses, and
memorability metadata, making it a natural substrate for studying a
brain-aligned memorability readout \citep{lahner2024boldmoments}. TRIBE v2 is a
multimodal brain-encoding model that maps video, audio, and language features to
predicted cortical responses \citep{tribev2}. In this paper, TRIBE is used as a
brain-aligned feature space, not as a literal brain scan or preference oracle.

\paragraph{Brain-to-media reconstruction.}
Recent brain decoding work reconstructs images or videos from measured fMRI
using diffusion, CLIP-style embeddings, semantic guidance, and temporal
reconstruction modules
\citep{lu2023minddiffuser,chen2023mindvideo,gong2024neuroclips,sun2024neurocine,yang2026semvideo}.
This is adjacent to our setting but distinct: those methods decode media from
brain activity, whereas our current selector scores generated media with
predicted-response and video-representation models.

\paragraph{Response-guided generation.}
Energy-guided diffusion uses a neuronal encoding model as an objective for
generating neurally exciting images \citep{pierzchlewicz2023energyguided}.
Saliency-guided generation methods similarly use predicted human attention to
condition generation or optimize diffusion latents
\citep{zhang2025gazefusion,wang2024sgool}. These papers motivate future
neural-response-guided video generation, but they do not turn our current
proxy-scored video candidates into human memorability evidence.

\paragraph{Video generation evaluation.}
Benchmarks such as VBench and FETV show that video generation evaluation is
multi-dimensional and that automatic metrics can diverge from human judgment
\citep{huang2023vbench,liu2023fetv}. Our selector therefore uses
memorability-specific scoring alongside preservation gates rather than replacing
quality and alignment checks.

\paragraph{Best-of-$N$ and preference alignment.}
Best-of-$N$ selection samples multiple candidates and returns the highest-scored
one under a reward model. Preference-optimization methods for image and video
diffusion models show a path toward amortizing such selectors into generators
\citep{wallace2023diffusiondpo,liu2024videodpo,wub2025densedpo}. We treat
distillation as future work because proxy-labeled DPO must still be validated
against humans.

\section{Learning a Brain-Aligned Memorability Direction}

\subsection{Dataset and features}

We analyze 1,022 clips from the 1,102-clip BOLD Moments dataset after feature and
URL exclusions. For each clip, we use saved TRIBE predicted cortical-response
features and mean-pool them across the available output time bins. We compare
against V-JEPA video embeddings and CLIP-style preservation scores when ranking
generated candidates.

\subsection{Readout}

Within each training fold, clips are sorted by BOLD Moments memorability. We
compute a unit contrastive readout using the top and bottom 30 percent of the
training set:
\[
v_{\mathrm{mem}} =
\frac{\mu_{\mathrm{high}}-\mu_{\mathrm{low}}}
{\lVert\mu_{\mathrm{high}}-\mu_{\mathrm{low}}\rVert_2}.
\]
Held-out clips are scored by projection onto $v_{\mathrm{mem}}$. This is a
supervised readout in a chosen representation frame, not a claim about the
ontological structure of memory.

\section{Current Results}

\subsection{BMD memorability prediction}

\begin{table}[h]
\centering
\caption{Current cross-validated memorability prediction on BOLD Moments.}
\begin{tabular}{llll}
\toprule
Predictor & Validation & Spearman $\rho$ & $n$ \\
\midrule
Gemini zero-shot & matched held-out subset & $0.139$ & 939 \\
V-JEPA contrastive & 5-fold CV & $0.395\pm0.037$ & 1,026 \\
TRIBE contrastive & 5-fold CV & $0.403\pm0.061$ & 1,022 \\
\bottomrule
\end{tabular}
\end{table}

The old headline that TRIBE cleanly beats V-JEPA is retired. The current result
is more interesting and more defensible: TRIBE is competitive with a strong
self-supervised video baseline while offering a brain-aligned substrate for
fMRI, ROI, and internal model-audit questions.

\subsection{Dominant-axis compactness}

The fold-safe linear ablation result supports dominant-axis compactness. Removing
the learned TRIBE direction and retraining on residual features drops
performance from roughly $\rho=0.401$ to $\rho=0.057$, while random direction
ablations leave the signal intact. Nonlinear residual probes recover some
signal, so the safe interpretation is dominant linear readout, not literal
one-dimensionality.

\subsection{Generated-video selector}

The current Wan2.2 proxy experiment evaluates 24 fresh prompt/image seeds. The
best product-style rule compares base clips to LoRA and best-of-$N$ candidates,
uses CLIP-style preservation gates, and keeps the base when the gated candidate
is worse. Under the TRIBE/BMD proxy, base-or-gated best-of-4 improves 18/24
seeds with mean lift $+2.817$ and median lift $+2.432$. This is a promising
selector-policy result, not independent behavioral evidence.

\begin{table}[h]
\centering
\caption{Wan2.2 proxy selector result. These are proxy scores, not human
validation.}
\begin{tabular}{llll}
\toprule
Selection rule & Improved seeds & Mean lift & Median lift \\
\midrule
Single LoRA & 20/24 & $+1.274$ & $+1.218$ \\
Raw best-of-$N$ & 20/24 & $+3.561$ & $+3.451$ \\
Base or raw best-of-$N$ & 20/24 & $+3.674$ & $+3.451$ \\
Base or gated best-of-$N$ & 18/24 & $+2.817$ & $+2.432$ \\
\bottomrule
\end{tabular}
\end{table}

\subsection{Representation-frame audit}

On the current candidate pool, TRIBE and V-JEPA are related but not identical.
Their scalar score Spearman correlation is $\rho=0.632$. Within
prompt-conditioned seed groups, their mean rank agreement is $\rho=0.443$, and
their top-1 selector agreement is 45.8 percent. This is the useful regime:
the baselines disagree often enough that human evaluation can adjudicate between
them.

\section{Planned Human Evaluation}

The decisive experiment is a blinded within-prompt pairwise study. Each trial
shows two videos generated for the same prompt or image seed. Participants choose
which clip they expect to be more memorable. The primary endpoint is the
prompt-clustered win rate of TRIBE+gate selected clips against the strongest
non-brain baseline, especially V-JEPA.

The current implementation assets include a V-JEPA-augmented manifest, pairwise
task file, and local Prolific-style survey HTML. No independent responses have
yet been collected for the V-JEPA-augmented selector pilot.

\section{Limitations}

First, the generated-video result is proxy-scored until human validation is
complete. Second, V-JEPA may tie or beat TRIBE in human judgments; if so, the
paper becomes a representation-frame comparison rather than a clean
brain-aligned-selector win. Third, memorability is not engagement, virality,
attention, quality, emotion, or commercial effectiveness. Fourth, direct
continuous steering is not solved; current evidence favors ranking and gating
candidate generations. Fifth, the TRIBE readout is a coordinate in a model frame,
not an ontological essence of memory.

\section{Conclusion}

Brain-aligned representations expose a compact memorability signal that is
competitive with strong video features and useful for defining a generated-video
selection workflow. The current result is strong enough to motivate a blinded
human validation study and careful IRB review. The submission-grade claim,
however, depends on that independent human endpoint.

\clearpage
\appendix
\section{Selector Human-Evaluation Protocol}

\subsection{Primary question}

When choosing among multiple generated videos for the same prompt, does a
TRIBE/BMD memorability selector choose clips that humans judge as more memorable
than clips chosen by non-brain baselines?

\subsection{Experimental unit}

The experimental unit is the prompt-conditioned candidate set. Each set contains
multiple videos generated from the same prompt or image seed. Selector policies
choose one candidate per set. Human raters compare videos only within the same
set, reducing confounding from prompt-level differences.

\subsection{Selector policies}

The minimum policy set is:
\begin{enumerate}
  \item random candidate selection;
  \item CLIP or text-video prompt-preservation winner;
  \item V-JEPA memorability winner;
  \item TRIBE/BMD memorability winner;
  \item TRIBE/BMD memorability winner with preservation or quality gate.
\end{enumerate}
Optional policies include aesthetic/video-quality winners, LoRA candidates, and
ensemble policies that combine memorability, quality, and prompt preservation.

\subsection{Human task}

The first study should use a fast pairwise preference task:

\begin{quote}
Which video do you think would be more memorable if you saw many such clips?
\end{quote}

This measures predicted memorability rather than delayed recognition. If this
pilot produces a clear effect, a follow-up delayed-recognition study should
present target and filler clips, wait for a delay, and then test recognition
against repeated and lure clips.

\subsection{Primary endpoint}

The primary endpoint is the prompt-clustered win rate of TRIBE+gate selected
clips against the strongest non-brain baseline. The primary hypothesis is:
\[
P(\text{human chooses TRIBE+gate over baseline}) > 0.5.
\]

\subsection{Statistical analysis}

The analysis should report:
\begin{itemize}
  \item prompt-clustered bootstrap confidence intervals;
  \item per-prompt win-rate distributions;
  \item mixed-effects logistic regression with selector policy as a fixed effect
        and random intercepts for prompt and, if identifiable, rater;
  \item multiple-comparison correction when comparing many selector policies.
\end{itemize}

\subsection{Anti-circularity rule}

No claim of human improvement can rely on TRIBE-scored gains alone. TRIBE may
train or select candidates, but the final endpoint must be independent human
judgment or delayed recognition.

\section{IRB-Relevant Study Summary}

The proposed first study is a minimal-risk online behavioral experiment. Adult
participants recruited through Prolific will view short generated videos and
make pairwise judgments about predicted memorability. The study does not collect
clinical data, biometric data, sensitive personal histories, or deception-based
responses. Responses will be analyzed in aggregate to compare selector policies.

The main foreseeable risks are boredom, fatigue, mild visual discomfort, and the
ordinary privacy risks associated with online survey data. Participants should be
warned that short videos may contain motion or rapid visual change, and they
should be allowed to stop at any time. The study should avoid graphic, sexual,
political persuasion, medical, or otherwise sensitive stimuli unless separately
reviewed.

Data should be stored without names, emails, IP addresses, or other direct
identifiers where possible. Prolific IDs, if collected for payment or duplicate
prevention, should be stored separately from response data or replaced with
pseudonymous participant IDs before analysis.


\bibliographystyle{unsrtnat}
\bibliography{references}

\end{document}
